% 图契约：正入射的辐射压力；几何、方向和公式与正文同一约定。
\documentclass[tikz,border=5pt]{standalone}
\usepackage{fontspec,xeCJK,amsmath,amssymb}
\IfFontExistsTF{Times New Roman}{\setmainfont{Times New Roman}}{\setmainfont{TeX Gyre Termes}}
\IfFontExistsTF{Heiti SC}{\setCJKmainfont[AutoFakeBold=1.5]{Heiti SC}}{\setCJKmainfont[AutoFakeBold=1.5]{FandolHei-Regular}}
\usetikzlibrary{arrows.meta,calc,patterns}
\definecolor{ink}{HTML}{243745}
\definecolor{blue}{HTML}{275DAB}
\definecolor{green}{HTML}{177C63}
\definecolor{red}{HTML}{BD4935}
\definecolor{orange}{HTML}{AA6C16}
\definecolor{purple}{HTML}{715393}
\begin{document}
\begin{tikzpicture}[x=1cm,y=1cm,font=\fontsize{9.5}{12}\selectfont,
v/.style={-{Latex[length=2mm,width=1.3mm]},line width=.85pt},
axis/.style={-{Latex[length=1.7mm,width=1.1mm]},ink!65,line width=.55pt},
aux/.style={ink!45,dashed,line width=.55pt},
lab/.style={fill=white,inner sep=1.8pt},
card/.style={draw=ink!25,fill=blue!3,rounded corners=2pt,text width=6.1cm,minimum height=.85cm,inner sep=6pt,align=center}]
\path[use as bounding box](-8,-3.6)rectangle(8,3.65);
\node[font=\bfseries\fontsize{10.5}{13}\selectfont,text width=15.1cm,align=center]at(0.00000,3.22000){M7-004　正入射的辐射压力};
\begin{scope}[shift={(-4,0)}]
\node[]at(-1.60000,1.80000){完全吸收};
\draw[ink,line width=1.5pt](-0.85000,-1.00000)--(-0.85000,1.00000);
\draw[v,blue](-2.40000,0.50000)--(-1.00000,0.50000);
\node[text=blue]at(-1.60000,0.95000){$I$};
\draw[v,red](-0.80000,0.00000)--(-0.30000,0.00000);
\node[]at(1.50000,1.80000){完全反射};
\draw[ink,line width=1.5pt](2.25000,-1.00000)--(2.25000,1.00000);
\draw[v,blue](0.70000,0.50000)--(2.10000,0.50000);
\node[text=blue]at(1.50000,0.95000){$I$};
\draw[v,red](2.30000,0.00000)--(2.80000,0.00000);
\draw[v,green](2.10000,-0.50000)--(0.70000,-0.50000);
\node[]at(0.00000,-1.70000){蓝：入射；绿：反射；红：压力方向};
\end{scope}
\node[card]at(4.00000,1.25000){$p_{\rm abs}=I/c=3.3\times10^{-6}\,\mathrm{Pa}$};
\node[card]at(4.00000,-1.05000){$p_{\rm refl}=2I/c=6.7\times10^{-6}\,\mathrm{Pa}$};
\node[text width=15.2cm,align=center]at(0.00000,-3.20000){正入射；完全吸收时不含反射箭头。};
\end{tikzpicture}\end{document}
